% Source: physical PDF pp. 150--155 (printed pp. 127--132), from the
% Section 27.3 heading through Eq. (27.3.24), immediately before Section 27.4.
% Inventory: Eqs. (27.3.1)--(27.3.24), five unnumbered display groups, and
% no citations, footnotes, figures, tables, or typographic dividers.
% The initial coverage manifest incorrectly ended the range at (27.3.23).
% Uncertainties: none.

\subsection{Gauge-Invariant Action for General Gauge Superfields}

Our experience in the previous section with supersymmetric Abelian gauge
theories suggests immediately that in a general non-Abelian gauge theory
the kinematic Lagrangian for the fields \(V_\mu^A(x)\), \(\lambda^A(x)\),
and \(D^A(x)\) should appear as part of the gauge-invariant generalization
of Eq.~\eqref{eq:27.2.6}:
\begin{equation}
  \mathcal{L}_{\mathrm{gauge}}
  =-\frac14\sum_A f_{A\mu\nu}f_A^{\mu\nu}
  -\frac12\sum_A\left(\bar\lambda_A\sl{D}\lambda_A\right)
  +\frac12\sum_A D_A D_A\,.
  \tag{27.3.1}\label{eq:27.3.1}
\end{equation}
We are now using a basis for the Lie algebra with totally antisymmetric
structure constants, and we are consequently not preserving the distinction
between upper and lower group indices, writing all indices \(A\), \(B\), etc.\
as subscripts. Also, \(f_{A\mu\nu}\) is the gauge-covariant field-strength
tensor
\begin{equation}
  f_{A\mu\nu}
  =\partial_\mu V_{A\nu}-\partial_\nu V_{A\mu}
  +\sum_{BC}C_{ABC}V_{B\mu}V_{C\nu}\,,
  \tag{27.3.2}\label{eq:27.3.2}
\end{equation}
and \(D_\mu\lambda\) is the gauge-covariant derivative of the gaugino field,
which in the adjoint representation is
\begin{equation}
  (D_\mu\lambda)_A
  =\partial_\mu\lambda_A+\sum_{BC}C_{ABC}V_{B\mu}\lambda_C\,.
  \tag{27.3.3}\label{eq:27.3.3}
\end{equation}
The question is: does Eq.~\eqref{eq:27.3.1} yield a supersymmetric action?

Since the Lagrangian density \eqref{eq:27.3.1} is manifestly
gauge-invariant, we can test whether the action is supersymmetric in any
convenient gauge. To find out whether
\(\delta\mathcal{L}_{\mathrm{gauge}}\) is a derivative at some point
\(X^\mu\), it is convenient to adopt a specific version of Wess--Zumino
gauge in which \(V_A^\mu(X)=0\). Then at \(X\) the changes in the component
fields are given by Eqs.~\eqref{eq:26.2.15}--\eqref{eq:26.2.17} at \(x=X\)
as
\begin{equation}
  \delta V_{A\mu}=(\bar\alpha\gamma_\mu\lambda_A)\,,
  \tag{27.3.4}\label{eq:27.3.4}
\end{equation}
\begin{equation}
  \delta\lambda_A
  =\left(
    \frac14 f_{A\mu\nu}[\gamma^\nu,\gamma^\mu]
    +i\gamma_5D_A
  \right)\alpha\,,
  \tag{27.3.5}\label{eq:27.3.5}
\end{equation}
\begin{equation}
  \delta D_A=i\left(\bar\alpha\gamma_5\sl{\partial}\lambda_A\right).
  \tag{27.3.6}\label{eq:27.3.6}
\end{equation}
(We must set \(x^\mu\) in these expressions equal to \(X^\mu\)
\textit{after} calculating the change under a supersymmetry transformation,
not before.) Also, the non-linear terms in \(f_A^{\mu\nu}\) are quadratic
in the \(V\)s and therefore at \(x=X\) have zero variation, so at \(x=X\)
\begin{equation}
  \delta f_{A\mu\nu}
  =\left(\bar\alpha
    (\gamma_\nu\partial_\mu-\gamma_\mu\partial_\nu)\lambda_A
  \right).
  \tag{27.3.7}\label{eq:27.3.7}
\end{equation}
With one exception, the terms in Eq.~\eqref{eq:27.3.1} and their
transformations under supersymmetry transformations are thus just a number
of copies (labelled by \(A\)) of the Abelian theory discussed in the
previous section, and therefore give a supersymmetric action. The one
exception that might disturb the supersymmetry of the action arises from
the second term in the gauge-covariant derivative \eqref{eq:27.3.3} of the
gaugino field:
\begin{equation}
  \mathcal{L}_{\lambda\lambda V}
  =-\frac12\sum_{ABC}C_{ABC}
    \left(\bar\lambda_A\sl{V}_B\lambda_C\right),
  \tag{27.3.8}\label{eq:27.3.8}
\end{equation}
whose variation at \(x=X\) is
\begin{equation}
  \begin{aligned}
  \delta\mathcal{L}_{\lambda\lambda V}
  &=-\frac12\sum_{ABC}C_{ABC}
    \left(\bar\lambda_A(\delta\sl{V}_B)\lambda_C\right)
  \\
  &=-\frac12\sum_{ABC}C_{ABC}
    \left(\bar\lambda_A\gamma_\mu\lambda_C\right)
    \left(\bar\alpha\gamma^\mu\lambda_B\right).
  \end{aligned}
  \tag{27.3.9}\label{eq:27.3.9}
\end{equation}
We can write the product of bilinears on the right-hand side as the sum
of two terms
\[
  \left(\bar\lambda_A\gamma_\mu\lambda_C\right)
  \left(\bar\alpha\gamma^\mu\lambda_B\right)
  =X_{ABC}+Y_{ABC}\,,
\]
with
\[
  \begin{aligned}
  X_{ABC}
  &\equiv\frac14\sum_{\pm}
    \left(\bar\lambda_A(1\mathbin{\pm}\gamma_5)\gamma_\mu\lambda_C\right)
    \left(\bar\alpha\gamma^\mu(1\mathbin{\pm}\gamma_5)\lambda_B\right),
  \\
  Y_{ABC}
  &\equiv\frac14\sum_{\pm}
    \left(\bar\lambda_A(1\mathbin{\pm}\gamma_5)\gamma_\mu\lambda_C\right)
    \left(\bar\alpha\gamma^\mu(1\mathbin{\mp}\gamma_5)\lambda_B\right).
  \end{aligned}
\]
By using standard Fierz identities and the anticommutativity of the spinor
fields, we have
\[
  \begin{aligned}
  &\left(\bar\lambda_A(1\mathbin{\pm}\gamma_5)\gamma_\mu\lambda_B\right)
   \left(\bar\alpha(1\mathbin{\pm}\gamma_5)\gamma^\mu\lambda_C\right)
  \\
  &\qquad=
   \left(\bar\lambda_A(1\mathbin{\pm}\gamma_5)\gamma_\mu\lambda_C\right)
   \left(\bar\alpha(1\mathbin{\pm}\gamma_5)\gamma^\mu\lambda_B\right),
  \\
  &\left(\bar\lambda_A(1\mathbin{\pm}\gamma_5)\gamma_\mu\lambda_B\right)
   \left(\bar\alpha(1\mathbin{\mp}\gamma_5)\gamma^\mu\lambda_C\right)
  \\
  &\qquad=
   \left(\bar\lambda_C(1\mathbin{\pm}\gamma_5)\gamma_\mu\lambda_B\right)
   \left(\bar\alpha(1\mathbin{\mp}\gamma_5)\gamma^\mu\lambda_A\right).
  \end{aligned}
\]
(To derive the first of these relations, we note that
\([(1\mathbin{\pm}\gamma_5)\gamma_\mu]_{\alpha\gamma}
 [(1\mathbin{\pm}\gamma_5)\gamma^\mu]_{\delta\beta}\)
may be thought of as the \(\alpha\beta\) matrix element of a matrix
depending on \(\delta\) and \(\gamma\), and may therefore be expanded in a
linear combination of \(\mathbf{1}_{\alpha\beta}\),
\((\gamma^\mu)_{\alpha\beta}\),
\([\gamma^\mu,\gamma^\kappa]_{\alpha\beta}\),
\((\gamma_5\gamma^\mu)_{\alpha\beta}\), and
\((\gamma_5)_{\alpha\beta}\). Because of the factors
\((1\mathbin{\pm}\gamma_5)\), the only term in the expansion is
proportional to
\([(1\mathbin{\pm}\gamma_5)\gamma^\mu]_{\alpha\beta}\). Lorentz invariance
and the presence of the other \(1\mathbin{\pm}\gamma_5\) factor tell us
that this expansion takes the form
\[
  [(1\mathbin{\pm}\gamma_5)\gamma_\mu]_{\alpha\gamma}
  [(1\mathbin{\pm}\gamma_5)\gamma^\mu]_{\delta\beta}
  =k[(1\mathbin{\pm}\gamma_5)\gamma_\mu]_{\alpha\beta}
    [(1\mathbin{\pm}\gamma_5)\gamma^\mu]_{\delta\gamma}.
\]
To determine the proportionality constant \(k\), we may contract both sides
with \((\gamma_\nu)_{\gamma\alpha}\) and find \(k=-1\). The minus sign is
cancelled by a minus sign arising from the anticommutation of \(\lambda_C\)
and \(\bar\alpha\). The other Fierz identity is proved in the same way,
except that we also need to use the symmetry property
\eqref{eq:26.A.7} for Majorana bilinears.) Hence \(X_{ABC}\) is symmetric
under interchange of \(B\) and \(C\), while \(Y_{ABC}\) is symmetric under
interchange of \(A\) and \(B\). Since \(C_{ABC}\) is totally
antisymmetric, both \(X_{ABC}\) and \(Y_{ABC}\) make a vanishing
contribution to the sum in Eq.~\eqref{eq:27.3.9}, leaving us
\(\delta\mathcal{L}_{\lambda\lambda V}=0\), so that
Eq.~\eqref{eq:27.3.1} gives a supersymmetric action, as was to be shown.

We can understand \textit{why} Eq.~\eqref{eq:27.3.1} gives a
supersymmetric action by identifying the superfield that has
\(f_{A\mu\nu}\), \(\lambda_A\), and \(D_A\) as component fields. Recall
that under a generalized gauge transformation, the vector superfield
\(V_A(x,\theta)\) has the transformation property \eqref{eq:27.1.12}:
\begin{equation}
  \begin{aligned}
  \exp\left(-2\sum_A t_A V_A(x,\theta)\right)
  \longrightarrow{}&
  \exp\left(-i\sum_A t_A\Omega_A(x,\theta)\right)
  \exp\left(-2\sum_A t_A V_A(x,\theta)\right)
  \\
  &{}\cdot
  \exp\left(+i\sum_A t_A\Omega_A^*(x,\theta)\right),
  \end{aligned}
  \tag{27.3.10}\label{eq:27.3.10}
\end{equation}
where \(\Omega_A(x,\theta)\) is a general left-chiral superfield. This is
not a gauge-covariant transformation rule, because
\(\Omega_A^*\ne\Omega_A\). To eliminate the factor involving
\(\Omega_A^*\), we note that \(\Omega_A^*\) is a right-chiral superfield,
so that \(\mathcal{D}_{L\alpha}\Omega_A^*=0\), and therefore
\begin{equation}
  \begin{aligned}
  &\exp\left(-2\sum_A t_A V_A(x,\theta)\right)
  \mathcal{D}_{L\alpha}
  \exp\left(+2\sum_A t_A V_A(x,\theta)\right)
  \\
  &\quad\longrightarrow
  \exp\left(-i\sum_A t_A\Omega_A(x,\theta)\right)
  \exp\left(-2\sum_A t_A V_A(x,\theta)\right)
  \\
  &\qquad\qquad{}\cdot
  \mathcal{D}_{L\alpha}\left[
    \exp\left(+2\sum_A t_A V_A(x,\theta)\right)
    \exp\left(+i\sum_A t_A\Omega_A(x,\theta)\right)
  \right].
  \end{aligned}
  \tag{27.3.11}\label{eq:27.3.11}
\end{equation}
\begin{sloppypar}
This is still not gauge-covariant, because the left-superderivative
\(\mathcal{D}_{L\alpha}\) acts on
\(\exp(+i\sum_A t_A\Omega_A(x,\theta))\) as well as on
\(\exp(+2\sum_A t_A V_A(x,\theta))\). This is eliminated if we follow
the lead of the Abelian theory discussed in the previous section, and
define a spinor superfield
\end{sloppypar}
\begin{equation}
  \begin{aligned}
  2\sum_A t_AW_{AL\alpha}(x,\theta)
  \equiv{}&
  \sum_{\beta\gamma}\epsilon_{\beta\gamma}
  \mathcal{D}_{R\beta}\mathcal{D}_{R\gamma}
  \left[
    \exp\left(-2\sum_A t_AV_A(x,\theta)\right)
  \right.
  \\
  &\left.\hspace{27mm}{}
    \cdot\mathcal{D}_{L\alpha}
    \exp\left(+2\sum_A t_AV_A(x,\theta)\right)
  \right].
  \end{aligned}
  \tag{27.3.12}\label{eq:27.3.12}
\end{equation}
Because the product of any three \(\mathcal{D}_R\)s vanishes,
\(W_{AL\alpha}\) is left-chiral
\begin{equation}
  \mathcal{D}_{R\beta}W_{AL\alpha}(x,\theta)=0\,,
  \tag{27.3.13}\label{eq:27.3.13}
\end{equation}
and because
\(\mathcal{D}_{R\beta}\mathcal{D}_{R\gamma}
\mathcal{D}_{L\alpha}\Omega_A
\propto\mathcal{D}_{R\delta}\Omega_A=0\),
\(W_{AL\alpha}\) is gauge-covariant in the sense that, for a generalized
gauge transformation,
\begin{equation}
  \begin{aligned}
  \sum_A t_AW_{AL\alpha}(x,\theta)
  \longrightarrow{}&
  \exp\left(-i\sum_A t_A\Omega_A(x,\theta)\right)
  \sum_A t_AW_{AL\alpha}(x,\theta)
  \\
  &{}\cdot
  \exp\left(+i\sum_A t_A\Omega_A(x,\theta)\right).
  \end{aligned}
  \tag{27.3.14}\label{eq:27.3.14}
\end{equation}

To calculate the spinor superfield at a point \(x^\mu=X^\mu\), we can
again adopt a version of Wess--Zumino gauge in which \(V_A(X)=0\), and
after a straightforward calculation find that in this gauge
\[
  \begin{aligned}
  W_{AL}(X,\theta)={}&
  \lambda_{AL}(X_+)
  +\frac12\gamma^\mu\gamma^\nu\theta_L
    \left(
      \partial_\mu V_{A\nu}(X_+)-\partial_\nu V_{A\mu}(X_+)
    \right)
  \\
  &+\left(\theta_L^{\mathsf T}\epsilon\theta_L\right)
    \sl{\partial}\lambda_{RA}(X_+)
  -i\theta_LD_A(X_+).
  \end{aligned}
\]
Since \(W_{AL}\) is gauge-covariant, at a general point in a general gauge
it must have the value
\begin{equation}
  \begin{aligned}
  W_{AL}(x,\theta)={}&
  \lambda_{AL}(x_+)
  +\frac12\gamma^\mu\gamma^\nu\theta_L f_{A\mu\nu}(x_+)
  +\left(\theta_L^{\mathsf T}\epsilon\theta_L\right)
    \sl{D}\lambda_{RA}(x_+)
  -i\theta_LD_A(x_+).
  \end{aligned}
  \tag{27.3.15}\label{eq:27.3.15}
\end{equation}
From this, we can construct a Lorentz- and gauge-invariant
\(\mathcal{F}\)-term bilinear in \(W\)
\begin{equation}
  \begin{aligned}
  -\left[
    \sum_{A\alpha\beta}
    \epsilon_{\alpha\beta}W_{AL\alpha}W_{AL\beta}
  \right]_{\mathcal{F}}
  =\sum_A\biggl[
    &-\left(\bar\lambda_A\sl{D}(1-\gamma_5)\lambda_A\right)
    -\frac12 f_{A\mu\nu}f_A^{\mu\nu}
  \\
    &+\frac{i}{4}\epsilon_{\mu\nu\rho\sigma}
      f_A^{\mu\nu}f_A^{\rho\sigma}
    +D_A^2
  \biggr].
  \end{aligned}
  \tag{27.3.16}\label{eq:27.3.16}
\end{equation}
Just as in the previous section, the gauge-invariant Lagrangian
\eqref{eq:27.3.1} is obtained from the real part of this
\(\mathcal{F}\)-term:
\begin{equation}
  -\frac12\operatorname{Re}\left[
    \sum_{A\alpha\beta}
    \epsilon_{\alpha\beta}W_{AL\alpha}W_{AL\beta}
  \right]_{\mathcal{F}}
  =\mathcal{L}_{\mathrm{gauge}}\,.
  \tag{27.3.17}\label{eq:27.3.17}
\end{equation}
What about the imaginary part? This is given by
\begin{equation}
  \begin{aligned}
  -\operatorname{Im}\left[
    \sum_{A\alpha\beta}
    \epsilon_{\alpha\beta}W_{AL\alpha}W_{AL\beta}
  \right]_{\mathcal{F}}
  ={}&
  -i\sum_A\left(\bar\lambda_A\sl{D}\gamma_5\lambda_A\right)
  \\
  &+\frac14\epsilon_{\mu\nu\rho\sigma}
    \sum_A f_A^{\mu\nu}f_A^{\rho\sigma}.
  \end{aligned}
  \tag{27.3.18}\label{eq:27.3.18}
\end{equation}
Eq.~\eqref{eq:26.A.7} and the antisymmetry of the structure constants
show that
\((\bar\lambda_A\sl{D}\gamma_5\lambda_A)
=\tfrac12\partial_\mu(\bar\lambda_A\gamma^\mu\gamma_5\lambda_A)\),
so the first term is a total derivative, while Eq.~\eqref{eq:23.5.4}
tells us that the second term is a total derivative also. In Abelian gauge
theories, this means that a term like \eqref{eq:27.3.18} would have no
effect. But as discussed in Sections 23.5 and 23.6, in non-Abelian gauge
theories the existence of instanton solutions allows the density
\eqref{eq:27.3.18} to have a non-vanishing integral over spacetime. We
therefore must allow for the possibility of a new term in the Lagrangian
density
\begin{equation}
  \mathcal{L}_\theta
  =-\frac{g^2\theta}{16\pi^2}
  \operatorname{Im}\left[
    \sum_{A\alpha\beta}
    \epsilon_{\alpha\beta}W_{AL\alpha}W_{AL\beta}
  \right]_{\mathcal{F}}\,,
  \tag{27.3.19}\label{eq:27.3.19}
\end{equation}
where \(\theta\) is a new real parameter, and \(g\) is a gauge coupling,
which can be conveniently defined for a simple gauge group so that if
\(t_A\), \(t_B\), and \(t_C\) are in the `standard' \(SU(2)\) subalgebra
of the gauge algebra used in calculating instanton effects, we have
\(C_{ABC}=g\epsilon_{ABC}\). With this definition of the gauge coupling,
Eq.~\eqref{eq:23.5.20} gives for simple gauge groups
\begin{equation}
  \int d^4x\,\epsilon_{\mu\nu\rho\sigma}
  \sum_A f_A^{\mu\nu}f_A^{\rho\sigma}
  =64\pi^2\nu/g^2\,,
  \tag{27.3.20}\label{eq:27.3.20}
\end{equation}
where \(\nu=0,\mathbin{\pm}1,\mathbin{\pm}2,\ldots\) is an integer, the
winding number, which characterizes the topological class of the gauge
field configuration. Thus for instantons of winding number \(\nu\), the
Lagrangian density \(\mathcal{L}_\theta\) contributes a phase to path
integrals, given by
\begin{equation}
  \left[
    \exp\left(i\int d^4x\,\mathcal{L}_\theta\right)
  \right]_\nu
  =\exp(i\nu\theta)\,,
  \tag{27.3.21}\label{eq:27.3.21}
\end{equation}
so the effects of \(\mathcal{L}_\theta\) are periodic in \(\theta\), with
period \(2\pi\).

It is often convenient to absorb a factor \(g\) into the gauge field, so
that the structure constants do not depend on \(g\), and instead the
Lagrangian density for the gauge field is multiplied by an over-all factor
\(1/g^2\). In this notation, the complete Lagrangian density for the gauge
field may be written in terms of the rescaled gauge fields and structure
constants as
\begin{equation}
  \mathcal{L}_{\mathrm{gauge}}+\mathcal{L}_\theta
  =-\operatorname{Re}\left[
    \frac{\tau}{8\pi i}
    \sum_{A\alpha\beta}
    \epsilon_{\alpha\beta}W_{AL\alpha}W_{AL\beta}
  \right]_{\mathcal{F}}\,,
  \tag{27.3.22}\label{eq:27.3.22}
\end{equation}
where \(\tau\) is the complex coupling parameter
\begin{equation}
  \tau=\frac{4\pi i}{g^2}+\frac{\theta}{2\pi}\,.
  \tag{27.3.23}\label{eq:27.3.23}
\end{equation}
According to Eq.~\eqref{eq:23.5.19}, the contribution of instantons of
winding number \(\nu\) to path integrals is suppressed by a factor
\(\exp(-8\pi^2|\nu|/g^2)\), which together with the factor
\eqref{eq:27.3.21} yields an over-all factor
\begin{equation}
  \exp\left[
    i\nu\theta-\frac{8\pi^2|\nu|}{g^2}
  \right]
  =
  \begin{cases}
    \exp(2\pi i\nu\tau), & \nu\geq 0,\\
    \exp(2\pi i\nu\tau^*), & \nu\leq 0.
  \end{cases}
  \tag{27.3.24}\label{eq:27.3.24}
\end{equation}
